% main.tex -- actual paper content.
% Template/formatting logic lives in intlpaper.cls; edit only content here.
% Compile with: xelatex main.tex && bibtex main && xelatex main.tex && xelatex main.tex

\documentclass{intlpaper}

\title{A Descriptive and Specific Title of Your Research Paper}

\paperauthors{%
  \authorblock{First A. Author}{Department of Computer Science\\University Name, City, Country}{first.author@university.edu}%
  \hfill
  \authorblock{Second B. Author}{Department of Computer Science\\University Name, City, Country}{second.author@university.edu}%
  \hfill
  \authorblock{Third C. Author}{Faculty of Engineering\\Another University, City, Country}{third.author@another.edu}%
}

\begin{document}

\maketitle

\begin{abstract}
This paper presents a concise summary of the research problem, the proposed
approach, and the main findings. The abstract should be self-contained,
typically 150--250 words, and state the motivation, method, key results, and
significance of the work without citing references or using undefined
abbreviations.
\end{abstract}

\begin{keywords}
keyword one, keyword two, keyword three, keyword four
\end{keywords}

\section{Introduction}
State the problem being addressed, why it matters, and how this work
advances the state of the art. Close the introduction with a short summary
of the paper's contributions and its organization, for example: ``The
remainder of this paper is organized as follows. Section~II reviews related
work. Section~III describes the proposed method. Section~IV presents the
experimental results. Section~V concludes the paper.''

A sample citation looks like this~\cite{smith2023}, and multiple citations
can be grouped together~\cite{lee2022,brown2021}.

\section{Related Work}
Summarize prior work relevant to the problem, positioning the contribution
of this paper against existing approaches~\cite{lee2022}.

\section{Methodology}
Describe the proposed method, dataset, or theoretical framework in enough
detail for the work to be reproducible.

\subsection{Problem Formulation}
Introduce notation and formal definitions here. Equations are numbered and
can be referenced later, e.g. Eq.~\eqref{eq:example}.

\begin{equation}
  y = f(x; \theta) = \sum_{i=1}^{n} \theta_i x_i + b
  \label{eq:example}
\end{equation}

\subsection{Proposed Approach}
Explain the core method, algorithm, or model architecture.

\section{Results and Discussion}
Present experimental results with supporting figures and tables.
Figure~\ref{fig:example} illustrates a placeholder result, and
Table~\ref{tab:example} summarizes example data.

\begin{figure}[t]
  \centering
  \fbox{\parbox[c][4cm][c]{0.9\linewidth}{\centering Figure placeholder --\\replace with \texttt{\textbackslash includegraphics}}}
  \caption{Example figure caption describing the content and key takeaway.}
  \label{fig:example}
\end{figure}

\begin{table}[t]
  \centering
  \caption{Example comparison of methods.}
  \label{tab:example}
  \begin{tabular}{lcc}
    \hline
    Method & Metric A & Metric B \\
    \hline
    Baseline    & 0.72 & 0.65 \\
    Proposed    & \textbf{0.85} & \textbf{0.79} \\
    \hline
  \end{tabular}
\end{table}

Discuss the results, comparing them against the baseline and related work,
and interpret their significance.

\section{Conclusion}
Summarize the main contributions, findings, and limitations, and suggest
directions for future work.

\section*{Acknowledgment}
Optional section acknowledging funding sources, institutions, or individuals
who contributed to the work but do not qualify as authors.

\bibliographystyle{ieeetr}
\bibliography{references}

\end{document}
